\addsection{Heroes}{\images/sorcery.png}
\hypertarget{Heroes}{Players} always control a main hero and may additionally also recruit a secondary hero. A “player’s hero” may refer to either of them. Heroes are used to perform movement actions on the game board and to start combats against enemies in order to reach a scenario victory condition.
\subsection*{Main hero}

The main hero is represented by its chosen model and hero card. Each faction’s main hero has 3 movement points. Only the main hero can use its own deck of cards.\par
Each main hero starts the game at level 1 and can advance up to level 7 by gaining experience. Experience is gained from winning combat, visiting certain locations and the treasure die. Gaining 1 experience is represented by the symbol \includesvg[height=10px]{\svgs/exp.svg}.

\subsection*{\hypertarget{Secondary}{Secondary heroes}}
If you control a town or a settlement, a secondary hero can be hired by flipping your population token and paying 10 gold. \textbf{Note}: Units cannot be \hyperlink{Units}{recruited or reinforced} during that use of the population token.\par
Your secondary hero uses the remaining hero model of your faction. You may wish to mark this model with a token such as a faction cube to differentiate it from the main hero. After hiring a secondary hero, place the model in a town or settlement you control. You can only have one secondary hero at a time.\par
Secondary heroes have 2 movement points. When you gain a secondary hero, take an additional set of 2 movement tokens to represent their MP. They do not have their own hero card, cannot gain experience, and cannot use your cards from your deck for any reason. If a secondary hero gains any cards, place them into your hand as normal.
Secondary heroes are considered to have the same level as the main hero for the purposes of resolving \hyperlink{Quick}{quick combat} against neutral enemies.\par
Units owned by the player may always participate in combat started in the same field as their secondary hero.\par
If your secondary hero is attacked by an enemy hero, you can choose to have your hero be \hyperlink{Endcombat}{instantly defeated instead of fighting a combat}. When a secondary hero is defeated, remove them from the game. They can be recruited again with another use of the population token.\par

\clearpage

\subsection*{\hypertarget{Herocard}{Hero Card Anatomy}}
\bigbreak
\begin{figure}[h]
  \begin{minipage}[t]{0.5\textwidth}
    \vspace{0pt}
    \begin{enumerate}[itemsep=5pt]
      \item Name – The hero’s name. Used for identification. Has no gameplay effect.
      \item Class – The hero’s class. Has no gameplay effect.
      \item Type – The hero’s type (might \includesvg[height=10px]{\svgs/might.svg} or magic \includesvg[height=10px]{\svgs/magic.svg}). Determines the amount of magic arrow spells in your starting deck (1 or 2 respectively).
      \item Faction color – Reminder for the color of the faction’s cubes and miniatures.
      \item Attack – Number of attack cards in your starting deck.
      \item Defense – Number of defense cards in your starting deck.
      \item Power – Number of power cards in your starting deck.
      \item Knowledge – Number of knowledge cards in your starting deck.
      \item Starting Ability – Reminder for the unique ability card the hero starts with.
      \item Hero Specialty – Reminder for the specialty cards the hero adds to their deck at the start of the game and after specific level ups. Each hero has three specialty cards.
      \item Level tracker – Whenever a main hero gains 1 or more experience \includesvg[height=10px]{\svgs/exp.svg} move the cube that number of steps on this track. When the cube reaches the next slot on the upper row, the hero gains a level.
    \end{enumerate}
  \end{minipage}\hfill
  \begin{minipage}[t]{0.48\textwidth}
    \centering
    \vspace{0pt}
    \begin{scriptsize}
      \begin{tikzpicture}
        \draw (0, 0) node[inner sep=0] {\makebox[\textwidth][c]{\includegraphics[width=\linewidth]{\images/herocard.png}}};
        \draw (2.3, 2.5) node {\encircle{\phantom{.}1\phantom{.}}};
        \draw (0.8, 1.9) node {\encircle{\phantom{.}2\phantom{.}}};
        \draw (3.5, 2.5) node {\encircle{\phantom{.}3\phantom{.}}};
        \draw (-0.1, 2.5) node {\encircle{\phantom{.}4\phantom{.}}};
        \draw (0, 1.25) node {\encircle{\phantom{.}5\phantom{.}}};
        \draw (1.1, 1.25) node {\encircle{\phantom{.}6\phantom{.}}};
        \draw (2, 1.25) node {\encircle{\phantom{.}7\phantom{.}}};
        \draw (3.25, 1.25) node {\encircle{\phantom{.}8\phantom{.}}};
        \draw (1, -0.2) node {\encircle{\phantom{.}9\phantom{.}}};
        \draw (3, -0.2) node {\encircle{10}};
        \draw (-1.7, -1.4) node {\encircle{11}};
      \end{tikzpicture}
    \end{scriptsize}
    \break
    \footnotesize{\textbf{\textit{\textcolor{purple}{Hero Card}}}}
    \scriptsize
    \begin{multicols}{2}
      \begin{itemize}
        \item[\textbf{1.}] \textbf{Name}
        \item[\textbf{2.}] \textbf{Class}
        \item[\textbf{3.}] \textbf{Type}
        \item[\textbf{4.}] \textbf{Faction color}
        \item[\textbf{5.}] \textbf{Attack}
        \item[\textbf{6.}] \textbf{Defense}
        \item[\textbf{7.}] \textbf{Power}
        \item[\textbf{8.}] \textbf{Knowledge}
        \item[\textbf{9.}] \textbf{Starting Ability}
        \item[\textbf{10.}] \textbf{Specialty}
        \item[\textbf{11.}] \textbf{Level tracker}
        \item[\textbf{\phantom{.}}] \phantom{.}
      \end{itemize}
    \end{multicols}
  \end{minipage}
\end{figure}

\clearpage

\subsection*{\hypertarget{Level}{Level effects}}
The most common sources of gaining experience are the treasure die and \hyperlink{Combatexperience}{combat}. Each new level up requires 2 experience. When a main hero reaches a new level, resolve the effects of the level up immediately. If a main hero gains multiple levels at the same time, resolve the effects in ascending level order. Gaining experience at level 7 has no effect.\par
The level tracker on your hero card shows the following information:
\begin{itemize}
\item Your main hero’s current level and amount of experience gained, shown by the cube's position.
\item Your current hand limit \includesvg[height=10px]{\svgs/hand.svg}.
\item The number of \hyperlink{Ability}{expert effects} \includesvg[height=10px]{\svgs/expert.svg} you may use during a round.
\item At which levels your main hero must search for a new \hyperlink{Ability}{ability card} or gain a \hyperlink{Specialty}{specialty card}. Level numbers written in gold on the level tracker (1, 4 and 6) give you a specialty card, while silver levels (2, 3, 5, 7) give you an ability card.
\end{itemize}
List of all effects:
\begin{itemize}
\item \textbf{Level 1} – Your hand limit is 4. Add your first specialty card to your deck.
\item \textbf{Level 2} – Search (2) the ability deck. You may play 1 card for their expert effect per round.
\item \textbf{Level 3} – Your hand limit is 5. Search (2) the ability deck.
\item \textbf{Level 4} – Gain your second specialty card. You may play 2 cards for their expert effect per round.
\item \textbf{Level 5} – Your hand limit is 6. Search (2) the ability deck.
\item \textbf{Level 6} – Gain your third specialty card. You may play 3 cards for their expert effect per round.
\item \textbf{Level 7} – Your hand limit is 7. Search (2) the ability deck.
\end{itemize}

\clearpage
\subsection*{\hypertarget{Playerdecks}{Player decks}}
All players have a unique deck which represents their main hero's abilities and equipment. Decks may contain statistic cards, ability cards, spell cards, artifact cards and the chosen hero’s specialty cards. Each player’s deck starts with 9 cards, built during the game’s set up. Game components may refer to this deck as your \textbf{Deck of Might and Magic}, in this rule book this is shorthanded to \textit{your deck}.
\subsection*{General card rules}
\begin{enumerate}
  \item Cards can be played \textbf{only on your turn}, or in a combat involving your \textbf{main hero}.
  \item After a card is used, discard it. Each player has their own separate discard pile for their own deck.
  \item If your deck is empty when you need to draw a card, shuffle your discard pile into a new deck to draw from.
  \item Whenever a hero gains a card for any reason, put it directly into your hand unless otherwise stated.
  \item Whenever you are instructed to \textbf{search} (X) the ability, artifact, or spell deck, you may either look at the top (X) cards from the specified deck, take one of them to your hand, and discard the others, or-instead of looking at the top (X) cards add the top card from that deck’s discard pile to your hand.
  \item The ability, artifact, and spell decks each have their own discard piles, created during the setup, which help players identify these decks. If a deck ever runs out of cards, reshuffle it and discard its top card to form a new discard pile. Also, whenever an ability, artifact or spell deck discard is empty, refill it with that deck’s top card.
  \item Player cards have the following types of effects:
  \begin{itemize}
        \item \textbf{Instant} \includesvg[height=10px]{\svgs/instant.svg} effects are resolved immediately and can be played out of turn during combat.
        \item \textbf{Activation} \includesvg[height=10px]{\svgs/activation.svg} effects must be played when activating your own unit in combat.
        \item \textbf{Map} \includesvg[height=10px]{\svgs/map.svg} effects cannot be used during combat.
        \item \textbf{Ongoing} \includesvg[height=10px]{\svgs/ongoing.svg} effects last until they are used up or until the player who played them starts their next turn (whichever happens first).
        \item \textbf{Permanent} \includesvg[height=10px]{\svgs/permanent.svg} cards stay in your play area until discarded or replaced. \textbf{Players may have only one permanent card at a time}; playing another discards the first.
    \end{itemize}

\end{enumerate}

\clearpage
\subsection*{\hypertarget{Ability}{Ability and statistic cards}}
\begin{wrapfigure}{R}{0.5\textwidth}

    \includegraphics[width=0.48\textwidth]{\images/ability.png}

\end{wrapfigure}

Every ability and statistic card has a basic effect and a stronger expert effect, which is shown below the basic effect. Whenever a player plays an ability or statistic card, they must choose which effect they are using. The number of expert \includesvg[height=10px]{\svgs/expert.svg} effects a player can use each round is limited by their main hero’s \hyperlink{Level}{level}. Track the number of uses each player has in any suitable manner, such as by moving black cubes on and off your hero card.\par
\textbf{Important}: Certain cards are limited to the Necropolis faction \includesvg[height=10px]{\svgs/necro.svg}. When a non-Necropolis player draws one, they may either discard it and draw a new card as a replacement or gain it as normal. Non-Necropolis players cannot use faction specific cards from their hand in any way besides discarding them.

\subsection*{Artifact cards}
Artifact cards are divided into 3 levels: minor, major, and relic. These levels relate to the overall power of the card and may be referenced when resolving certain effects or during scenario set up. Artifacts are gained through map exploration.\par
Artifacts can be \hyperlink{Trading}{traded} in alliance and cooperative scenarios.\par
\includegraphics[width=0.48\textwidth]{\images/artifact.png}


\clearpage

\subsection*{\hypertarget{spells}{Spell cards}}
\begin{wrapfigure}{R}{0.5\textwidth}
    \begin{center}
    \includegraphics[width=0.48\textwidth]{\images/spell_example.png}
    \end{center}
\end{wrapfigure}

Spell cards have three possible primary effects. Using the topmost basic version of the spell has no additional costs. Accessing the second and third effects requires the spell to be \textbf{empowered} before it is played. The amount a spell needs to be empowered for each effect is indicated within a circle in the middle of its card text.\par
You may empower a spell by playing any card for its empower \includesvg[height=10px]{\svgs/empower.svg} effect. The bottom, alternative effect of all spell cards is to empower another spell.\par
\textbf{Important}: During combat, \textbf{only one spell} card may be played by each player \textbf{per combat round}.\par
Spells can be gained by building the mage guild. When you do, search (2) the spell deck, twice. If you start a scenario with the mage guild, perform the searching in turn order and shuffle the cards into your deck at the end of the set up. More spells can be bought by flipping your spell book token and paying the indicated cost on the guild to search (2) the spell deck. This cannot be done during the same round as when the guild is constructed. Spells can be \hyperlink{Trading}{traded} in alliance and cooperative scenarios.

\subsection*{\hypertarget{Specialty}{Hero specialty cards}}
Hero specialty cards are gained from level ups. Each hero has an unique set of them. While many specialty cards have effects which resemble spell cards, specialty cards are their \textbf{own unique category of card}, which are not affected by effects that do not specify them.
\par

\begin{figure}[h]
\centering
\includegraphics[width=0.25\linewidth]{\images/specialty.png}
\end{figure}
\begin{center}
\textit{A level 1 specialty card, belonging to Rion the cleric.}
\end{center}

\clearpage
